\section{Introduction}\label{intro}

Large language models (LLMs) have emerged as powerful tools for code generation, capable of translating natural language specifications into executable programs. Recent advances demonstrate that LLMs can design algorithms that outperform hand-crafted heuristics~\cite{RomeraParedes2024FunSearch,Zhang2025SolSearch,Liu2024EoHS}. Yet a fundamental limitation persists: these systems rely on manual human intervention to interpret results, diagnose failures, and formulate refined prompts. For example, FunSearch~\cite{RomeraParedes2024FunSearch} requires researchers to manually design the evolutionary scaffold, select which programs to seed the initial population, and interpret whether discovered functions constitute genuine improvements; EvoPrompt~\cite{Chen2024EvoPrompt} relies on human-crafted meta-prompts that encode domain-specific guidance without automated failure diagnosis; and SolSearch~\cite{Zhang2025SolSearch} depends on human researchers to select benchmark suites, interpret PAR-2 score trends, and decide when to terminate optimization. The human acts as the critical feedback loop, transforming LLM collaboration from an automated discovery process into an assisted discovery process.

This study addresses this limitation by proposing SolEvolve, a framework that closes the feedback loop, transforming LLMs from passive code generators into active components within an autonomous evolutionary system. Where existing approaches require humans to manually analyze solver outputs, detect stagnation patterns, and craft new prompts, SolEvolve automates these steps through a four-component architecture: Generator (LLM-driven code synthesis), Evolver (population-based search), Verifier (rigorous validation), and Reflector (automated feedback synthesis). The resulting framework addresses a broader challenge than any single predecessor: fully automated algorithm discovery for multi-modal combinatorial search, with mathematically rigorous verification and closed-loop semantic adaptation.

We validate SolEvolve through progressively sophisticated case studies in combinatorial mathematics, where the exponential search space and rigorous correctness requirements provide ideal testbeds for automated discovery. The search for optimal error-correcting codes exemplifies the challenges SolEvolve addresses. For an $[n,k,d]$ binary linear code, the parity portion of the generator matrix contributes $k(n-k)$ Boolean variables, and the minimum-distance constraint introduces highly non-linear relationships among $2^{k}-1$ non-zero codewords~\cite{Huffman}. Traditional approaches rely on manual SAT encoding design or bespoke heuristics, demanding extensive domain expertise. Error-correcting codes provide an ideal validation domain because algebraic bounds (Singleton, Hamming, Gilbert--Varshamov) offer tight necessary or sufficient conditions~\cite{Huffman,Kim:2023}, enabling rigorous verification of discovered solutions. Beyond binary codes, we validate SolEvolve's generality through two additional domains: generalized Lucas cubes, where ILP formulations discover perfect partitions below Mollard's boundary~\cite{Mollard2022}; and self-orthogonal embeddings, where the Generator conceived a provably optimal rank-one elimination algorithm. This diversity demonstrates that SolEvolve's architecture generalizes across combinatorial domains, not merely codes.

We organize our contributions around the following capabilities of the SolEvolve framework:
First, we introduce an LLM-driven meta-optimization architecture that dynamically orchestrates symbolic constraint solvers (SAT/ILP) and heuristic methods (GA). The Reflector agent provides semantic diagnostics that guide the Generator's strategy, enabling adaptive problem partitioning.
Second, we present a SAT-GA hybrid approach that leverages verified mathematical constraints to repair structurally stagnating heuristic populations. Ablation studies confirm that this semantic feedback loop increases candidate upgrade success rates from 2\% (fitness-only baseline) to 10\%.
Third, we demonstrate generality across combinatorial domains: generalized Lucas cubes (first perfect partition of $\Lambda_7(1^4)$ below Mollard's boundary), and self-orthogonal embeddings across $\mathbb{F}_2, \mathbb{F}_3, \mathbb{F}_4, \mathbb{F}_5$.
Fourth, we provide a detailed analysis of LLM contamination. By introducing a structurally novel problem variant (zero-length embedding), we demonstrate the Generator's capacity for genuine algorithm synthesis rather than mere retrieval from pre-training data.

The remainder of this paper is organized as follows. Section~\ref{sec:related} reviews related work on LLM-evolutionary algorithm integration and automated heuristic design. Section~\ref{ch:design} presents the SolEvolve framework architecture, including the Generator-Evolver-Verifier-Reflector components and the hybrid SAT-GA methodology that bridges exact and heuristic methods. Sections~\ref{ch:codes}--\ref{ch:lucas} in the Appendix demonstrate applications across three combinatorial domains: binary linear codes, self-orthogonal embeddings, and generalized Lucas cubes. Finally, Section~\ref{ch:discussion} discusses implications and future directions, followed by conclusions in Section~\ref{ch:conclusion}.
